\textbf{پرسش ۱. مفاهیم نظری (۴۰ نمره)}

\begin{itemize}
    \item[\textbf{(الف)}] ثابت کنید یا نشان دهید که گزاره نادرست است: یک MLP با یک لایه پنهان و بدون تابع فعال‌سازی می‌تواند کمینه‌های محلی‌ای داشته باشد که کمینه سراسری نباشند.
    
    \item[\textbf{(ب)}] در شبکه‌های عمیق، با وجود غیرمحدب بودن تابع هزینه، روش‌های گرادیانی اغلب در عمل موفق‌اند. این موضوع را با تکیه بر هندسه نقاط بحرانی در فضاهای پربعد و طیف هسین توضیح دهید.
    
    \item[\textbf{(پ)}] به‌روزرسانی نزول گرادیان به صورت زیر است:
    \[
    \theta_{t+1} = \theta_t - \eta \nabla L(\theta_t).
    \]
    فرض کنید در گام $t$، گرادیان برابر باشد با
    \[
    \nabla L = \begin{bmatrix} 100 \\ 0.01 \end{bmatrix}.
    \]
    \begin{enumerate}
        \item[(۱)] توضیح دهید که چرا استفاده از یک نرخ یادگیری ثابت $\eta$ باعث ایجاد یک دوگانگی برای همگرایی می‌شود.
        \item[(۲)] به طور خلاصه توضیح دهید که روش‌های نرخ یادگیری تطبیقی مانند Adam یا RMSProp از نظر ریاضی چگونه این مشکل را حل می‌کنند. این روش‌ها بر چه کمیتی نرمال‌سازی انجام می‌دهند؟
    \end{enumerate}
    
    \item[\textbf{(ت)}] یک شبکه عصبی با تعداد پارامترهای بسیار بیشتر از تعداد نمونه‌ها را در نظر بگیرید، یعنی
    \[
    p \gg n
    \]
    و فرض کنید این شبکه روی $n$ نمونه آموزش داده می‌شود تا به خطای آموزش صفر برسد. اکنون فرض کنید \%۱۰ از برچسب‌های آموزشی به صورت تصادفی برعکس شوند و شبکه دوباره آموزش داده شود.
    \begin{enumerate}
        \item[(۱)] توضیح دهید چرا این شبکه بیش‌پارامتری‌شده، با وجود نویز در برچسب‌ها، هنوز می‌تواند به خطای آموزش صفر برسد.
        \item[(۲)] با وجود دقت کامل روی دادهٔ آموزش، خطای آزمون به طور قابل توجهی افزایش می‌یابد. این موضوع را با استفاده از مفهوم برازش سیگنال در برابر نویز و این که بیش‌پارامتری‌بودن چگونه بر چیزی که مدل یاد می‌گیرد اثر می‌گذارد، توضیح دهید.
    \end{enumerate}
\end{itemize}

\vspace{1cm}

\textbf{پرسش ۲. پس‌انتشار و عبور گرادیان (۳۵ نمره)}

یک MLP با ساختار زیر را در نظر بگیرید:
\begin{align*}
    &\text{ورودی:} & &d \\
    &\text{لایه پنهان اول:} & &h_1 \text{ نورون با تابع فعال‌سازی } \text{ReLU} \\
    &\text{لایه پنهان دوم:} & &h_2 \text{ نورون با تابع فعال‌سازی } \text{sigmoid} \\
    &\text{لایه خروجی:} & &K \text{ نورون با تابع فعال‌سازی } \text{softmax}
\end{align*}
تابع هزینه نیز cross-entropy است.

\begin{itemize}
    \item[\textbf{(الف)}] مشتق
    \[
    \frac{\partial L}{\partial W^{(1)}}
    \]
    را با استفاده از پس‌انتشار به دست آورید. همه مراحل قاعده زنجیره‌ای را نشان دهید و به نامشتق‌پذیر بودن ReLU در صفر اشاره کنید.
    
    \item[\textbf{(ب)}] فرض کنید همه فعال‌سازی‌های لایه پنهان دوم در ناحیه اشباع باشند، یعنی
    \[
    a_i^{(2)} \in \{0.01, 0.99\}.
    \]
    بیشینه قدر مطلق
    \[
    \frac{\partial L}{\partial W_{ij}^{(1)}}
    \]
    را بر حسب $\|W^{(3)}\|$، $\|W^{(2)}\|$ و $\|\delta^{(3)}\|$ به دست آورید. نشان دهید که چگونه رابطه
    \[
    \sigma'(z) = \sigma(z)(1 - \sigma(z))
    \]
    باعث می‌شود گرادیان تقریباً دو مرتبه بزرگی کوچک شود.
\end{itemize}

\vspace{1cm}

\textbf{پرسش ۳. مبانی نظری بهینه‌سازی (۲۵ نمره)}

برای نزول گرادیان
\[
\theta_{t+1} = \theta_t - \eta \nabla L(\theta_t)
\]
داریم:
\begin{enumerate}
    \item[(۱)] بسط تیلور مرتبه دوم $L(\theta_{t+1})$ را حول $\theta_t$ بنویسید.
    \item[(۲)] شرطی روی $\eta$ به دست آورید که کاهش تابع هزینه را تضمین کند، در حالتی که بیشینه مقدار ویژه هسین برابر $\lambda_{\max}$ است.
\end{enumerate}